% -*- coding: utf-8 -*-

\chapter{Introduction}

\section{The Role of Logs in Modern Operations}

Logs are the most pervasive and continuously available telemetry in modern software systems.
Unlike metrics, which aggregate numerical measurements, or distributed traces, which require
explicit instrumentation at the call level, logs are emitted by virtually every system component
as a natural by-product of execution~\cite{he2020loghub}. They record timestamped events,
exception messages, state transitions, and resource identifiers in semi-structured natural
language, making them the primary artifact that on-call engineers consult when a failure occurs.
In cloud-native architectures that span microservices, container orchestration platforms, and
distributed storage systems, a single production incident may leave traces across dozens of
components simultaneously, each producing its own log stream. Correlating and interpreting these
streams within the tight time windows demanded by service-level objectives (SLOs) is a task that
has grown in difficulty commensurate with the growth of system scale and
complexity~\cite{zhang2024cloudRCAsurvey}.

The practical importance of log analysis has motivated a substantial body of research in automated
log parsing, anomaly detection, and failure diagnosis. Log parsers such as Drain convert raw,
free-form log messages into structured event templates by masking variable
tokens~\cite{he2017drain,zhu2019tools}. These templates serve as the input to downstream analytical
pipelines, enabling quantitative comparison of log behaviors across time windows. However, log
analysis and root cause analysis (RCA) are distinct tasks with different objectives: while log
analysis broadly encompasses anomaly detection, classification, and summarization, RCA is
specifically concerned with identifying the underlying causal fault that produced an observable
incident and providing an actionable explanation for operators. This thesis is concerned
exclusively with the RCA task.

\section{Root Cause Analysis: Definition and Challenges}

Root cause analysis is the process of tracing an observed failure back to its primary causal
factor, as distinct from its immediate symptoms. In a microservice system, for example, a
database connection timeout may surface as an HTTP 500 error in an upstream API; the symptom is
the HTTP error, but the root cause is the database unavailability. Correct identification of the
root cause is prerequisite to effective remediation: treating the symptom (e.g., restarting the
upstream service) without addressing the root cause (e.g., the exhausted database connection pool)
will not prevent recurrence~\cite{meng2020localizing,wu2020microrca}.

The RCA task in production environments faces several well-documented challenges that are
difficult to address simultaneously. First, production systems emit logs at a rate that makes
manual inspection infeasible. A single Hadoop cluster, for example, can emit hundreds of thousands
of log lines per job execution, as evidenced by the Hadoop1 benchmark used in this
thesis~\cite{he2020loghub}. Engineers cannot review this volume manually during an active incident.
Second, failures in distributed systems rarely confine themselves to a single component. Fault
propagation across service dependencies means that the failure signal in the logs of the faulty
component may be temporally and spatially separated from the logs of the component that originally
triggered the cascade~\cite{wang2018cloudranger,chen2014causeinfer}. Third, the heterogeneity
of log formats across different software stacks, versions, and configurations means that pattern
matching rules and templates crafted for one system are rarely transferable to another without
substantial manual effort~\cite{zhang2019logrobust}. Fourth, logs encode domain-specific
terminology and operational conventions that require contextual knowledge to interpret correctly.
An error string like ``BlockInfo not found in volumeMap'' is diagnostic only to an engineer who
knows the semantics of HDFS block management. Capturing and operationalizing this domain knowledge
at scale is a fundamental bottleneck in automated RCA.

Traditional automated RCA approaches address these challenges through two broad strategies.
The first class of approaches employs statistical correlation and causal graph methods: they
construct dependency graphs between services or metrics and identify components whose
behavior changes are most causally associated with the observed
anomaly~\cite{wang2018cloudranger,meng2020localizing,chen2014causeinfer}. While effective in
controlled environments where component relationships are well-modeled, these methods require
significant upfront engineering effort to build and maintain the causal model, and they are
brittle when system topology changes frequently. The second class of approaches applies supervised
machine learning to structured log representations: log templates are extracted, converted into
feature vectors or event sequences, and fed into classifiers or sequence models trained to
recognize fault patterns~\cite{xu2009consolelogs}. These methods can achieve strong performance
on benchmark datasets, but they depend on labeled training data from the target system, are
sensitive to log statement drift caused by software updates~\cite{zhang2019logrobust}, and
produce predictions without interpretable explanations---a significant drawback when the goal is
to support operator decision-making rather than merely raise an alert.

\section{LLM-Based Attempts at RCA: Capabilities and Limitations}

The emergence of large language models (LLMs) with strong instruction-following and natural
language generation capabilities~\cite{brown2020languagemodelsfewshotlearners,ouyang2022traininglanguagemodelsfollow}
has motivated a new generation of RCA approaches that treat the diagnosis task as a prompted
reasoning problem. The central appeal is clear: an LLM can read heterogeneous log text without
requiring a domain-specific parser, generate a human-readable explanation of the likely root cause,
and suggest remediation steps---all in a zero-shot or few-shot setting that requires no task-specific
training data.

Several recent systems demonstrate the viability of this direction. RCACopilot, developed at
Microsoft and evaluated on a year's worth of production cloud incidents, employs an LLM to
aggregate diagnostic information collected by per-alert handlers and predict the root cause
category, achieving up to 76.6\% accuracy on real incident
data~\cite{chen2024rcacopilot}. Ahmed et al.\ conduct the first large-scale study of GPT-3.x
models for recommending root causes and mitigation steps from incident reports at Microsoft,
finding that LLMs can surface relevant diagnostic context even without task-specific
fine-tuning~\cite{ahmed2023rcarecommend}. Roy et al.\ explore the use of LLM-based
ReAct~\cite{yao2023reactsynergizingreasoningacting} agents equipped with retrieval tools for RCA
on an out-of-distribution dataset of production incidents, showing that agentic tool use can
improve diagnostic coverage but also identifying key failure modes when retrieval returns
irrelevant context~\cite{roy2024llmagentRCA}. AetherLog, which is the most directly related
prior work to this thesis, integrates an LLM with a knowledge graph of historical incidents for
log-based RCA, demonstrating that structured retrieval of similar past cases substantially
reduces unsupported speculation and improves the precision of generated
diagnoses~\cite{cui2025aetherlog}.

Despite these promising results, single-LLM and single-agent RCA approaches share a set of
structural limitations that have not been fully resolved. The most significant is the tendency
toward hallucination: when the prompt does not contain sufficient or unambiguous evidence, an LLM
will produce a plausible-sounding but unsupported diagnosis, and there is no internal mechanism
to detect or correct this~\cite{du2023improvingfactualityreasoninglanguage}. This is particularly
consequential in RCA because an incorrect diagnosis can lead an engineer to perform unnecessary
or harmful remediation actions. A related limitation is the absence of hypothesis diversity: a
single model pass commits to one explanation without systematically considering alternatives. In
cases where evidence is ambiguous---which is common when failures propagate across multiple
components---premature commitment to a single hypothesis risks anchoring the operator on the
wrong causal narrative. Furthermore, single-pass generation provides no calibrated confidence
signal: the model's expressed confidence correlates poorly with actual correctness, making it
difficult for operators to decide when to trust the automated diagnosis and when to escalate
to manual investigation~\cite{chen2024rcacopilot,roy2024llmagentRCA}.

\section{Addressing the Gap: Agent Debate with Knowledge Graph Grounding}

This thesis proposes \systemname{}, a multi-agent, knowledge-graph-guided RCA system designed
to address the limitations of single-agent LLM approaches. The core insight is that the
reliability problems of single-agent RCA stem from two compounding factors: insufficient
evidence grounding and the absence of adversarial verification. We address both simultaneously.

Evidence grounding is improved by integrating a Neo4j-based knowledge graph (KG) that stores
historical incidents, their involved entities, and their diagnosed root causes. At diagnosis time,
the KG is queried for historical incidents that share entities with the current incident, and the
retrieved context is provided to the reasoning agents as structured external memory. This design
is motivated by the success of retrieval-augmented generation for knowledge-intensive
tasks~\cite{lewis2021retrievalaugmentedgenerationknowledgeintensivenlp} and by the specific
observation that operational incident history---when organized as a structured
graph~\cite{hogan2021knowledge}---provides precisely the kind of entity-centric, relation-aware
context that supports reliable causal inference.

Adversarial verification is implemented through a structured debate protocol in which three
specialized reasoning agents---a log-focused agent, a KG-focused agent, and a hybrid
agent---independently generate competing hypotheses. A dedicated judge agent then evaluates all
hypotheses against a fixed rubric covering evidence quality, reasoning strength, confidence
calibration, completeness, and consistency. This design is motivated by recent empirical evidence
that multi-agent debate substantially improves the factual accuracy of LLM-generated
content~\cite{du2023improvingfactualityreasoninglanguage}, and by the observation that
specialization and decomposition reduce the chance that all agents share the same blind
spots~\cite{wu2023autogen}. By forcing hypotheses to compete under an explicit scoring
mechanism, the system reduces the risk of majority-class collapse and overconfident but
unsupported diagnoses that are the primary failure modes of single-agent baselines.

\section{Research Questions}

The following research questions structure the experimental investigation in this thesis:

\begin{description}
  \item[RQ1:] Does a multi-agent debate approach achieve higher diagnostic accuracy and better
              class balance than single-agent LLM baselines for log-based RCA?
  \item[RQ2:] Does structured debate with explicit judge scoring reduce overconfident or
              unsupported predictions compared to single-pass generation?
  \item[RQ3:] Does knowledge graph integration, as an external grounding source, measurably
              improve diagnosis quality beyond what retrieval-augmented single-agent generation
              achieves?
  \item[RQ4:] Do the generated explanations provide sufficient evidence citation and causal
              reasoning to support operator triage?
\end{description}

\section{Contributions}

This thesis makes the following contributions:

\paragraph{Multi-agent RCA architecture.}
We design and implement \systemname{}, a six-agent RCA system comprising a Log Parser, a KG
Retrieval Agent, three specialized Reasoners (log-focused, KG-focused, and hybrid), and a Judge
Agent. Each agent has a well-defined role and operates over structured input/output contracts,
enabling systematic composition and debugging.

\paragraph{Structured debate protocol.}
We propose a multi-round debate protocol in which three reasoning agents independently generate
competing hypotheses, a judge evaluates them against a fixed five-criterion rubric, and agents
refine their hypotheses based on judge feedback in subsequent rounds. This protocol operationalizes
reliability as a measurable, evidence-based selection process rather than an informal heuristic.

\paragraph{Knowledge graph for operational memory.}
We implement a Neo4j-based incident knowledge graph that supports entity-centric retrieval of
similar historical incidents. The KG is populated incrementally during evaluation, simulating
an operational deployment where diagnostic history accumulates over time.

\paragraph{Empirical evaluation across three datasets.}
We evaluate \systemname{} on three publicly available log datasets spanning two domains:
Hadoop1 and HDFS\_v1 from LogHub~\cite{he2020loghub}, and the CMCC OpenStack failure
dataset~\cite{cui2025aetherlog}. Across all three datasets, the multi-agent system consistently
outperforms single-agent and RAG baselines in macro-averaged \fone{} score, with the most
substantial improvement on CMCC, where single-agent baselines collapse to uninformative
predictions.

\paragraph{Fully local and reproducible implementation.}
The entire system runs on local infrastructure using Ollama for LLM inference and Neo4j Community
Edition for graph storage, with no dependence on commercial API availability. All prompts,
intermediate outputs, and evaluation scripts are logged, enabling full reproduction of results.

\section{Thesis Organization}

Chapter~2 reviews the background and related work, covering traditional and ML-based RCA,
knowledge graphs for operational reasoning, LLM foundations from transformers to instruction
tuning and alignment, and multi-agent reasoning systems. Chapter~3 presents the system design
and architecture of \systemname{}, detailing the six agents, the knowledge graph schema, and the
debate protocol. Chapter~4 describes the experimental methodology, including label normalization,
baseline configurations, and evaluation metrics. Chapter~5 covers the implementation, including
local LLM inference, structured output handling, and the knowledge graph pipeline.
Chapter~6 presents the experimental results on all three datasets with detailed per-class
analysis and a misclassification audit. Chapter~7 discusses the findings, limitations, and threats
to validity. Chapter~8 concludes with a summary of contributions, answers to the research
questions, and directions for future work. Appendix~A provides representative case studies
extracted from system outputs.
